\documentclass[12pt,a4paper]{article}
\usepackage[utf8]{inputenc}
\usepackage{graphicx}
\usepackage{float}
\usepackage{amsmath}

\title{Informe Taller 3}
\author{Jose Salazar, Sofia Sepulveda}
\date{\today}

\begin{document}
\maketitle

\section{Introducción}
El algoritmo \textit{Anytime Weighted A*} (AW-A*) es una extensión de A* que permite obtener una solución rápidamente y mejorarla de manera progresiva en el tiempo. 

\section{Resultados Experimentales}
\subsection{Pregunta A}
Disminuir el valor de $w$ en la función $f(n) = g(n) + w \cdot h(n)$ reduce la sobreestimación de la heurística. 
Con $w$ alto, AW-A* entrega soluciones rápidas pero subóptimas; al acercarse $w \to 1$, el algoritmo se comporta como A*, encontrando soluciones más cercanas a la óptima.
\subsection{Pregunta B}
En los mapas de $41 \times 41$, la primera solución encontrada es la solucion de optimalidad.
Al ser un mapa tan pequeño, el algoritmo no necesita mucho tiempo para encontrar la solución de optimalidad.

\subsection{Pregunta C}

\begin{table}[H]
    \centering
    \resizebox{\textwidth}{!}{%
    \begin{tabular}{|c|c|c|c|c|c|c|c|c|}
    \hline
    \textbf{Size} & \textbf{Manhattan} & \textbf{w\_start} & \textbf{Solved Rate} & \textbf{Avg. BSL} & \textbf{Avg. BST} & \textbf{Avg. FST} & \textbf{Avg. FSL} \\
    \hline
    (1500,1300) & False & 2.0 & 0.5 & 12827.0 & 35.2 & 10.2 & 12875.0 \\
    (1500,1300) & False & 2.5 & 0.5 & 12843.0 & 40.9 &  9.6 & 12875.0 \\
    (1500,1300) & False & 3.0 & 0.5 & 12843.0 & 51.1 &  8.9 & 12875.0 \\
    (1500,1300) & True  & 2.0 & 0.5 & 12827.0 & 34.9 & 10.2 & 12875.0 \\
    (1500,1300) & True  & 2.5 & 0.5 & 12843.0 & 42.6 &  9.7 & 12875.0 \\
    (1500,1300) & True  & 3.0 & 0.5 & 12843.0 & 51.0 &  8.8 & 12875.0 \\
    (2000,2000) & False & 2.0 & 0.5 & 16739.7 & 52.1 & 22.8 & 16751.0 \\
    (2000,2000) & False & 2.5 & 0.5 & 16744.3 & 55.7 & 20.2 & 16751.0 \\
    (2000,2000) & False & 3.0 & 0.5 & 16796.3 & 59.9 & 16.6 & 16897.0 \\
    (2000,2000) & True  & 2.0 & 0.5 & 16739.7 & 52.1 & 22.8 & 16751.0 \\
    (2000,2000) & True  & 2.5 & 0.5 & 16744.3 & 55.3 & 20.2 & 16751.0 \\
    (2000,2000) & True  & 3.0 & 0.5 & 16796.3 & 59.5 & 16.4 & 16897.0 \\
    \hline
    \end{tabular}}
    \caption{Resultados en mapas grandes: tasa de éxito, longitudes promedio y métricas de tiempos.}
    \label{tab:mapas_grandes_full}
\end{table}

Las pruebas con distintas combinaciones de $w$ y heurísticas muestran que la heurística Manhattan, junto con $w_{start} = 2.5$, ofrece una buena solucion en un tiempo mas rapido y la evolucion con el tiempo retorna soluciones de una buena calidad. 
La heurística Euclidiana no aporta mejoras significativas considerando su mayor costo computacional. 
La tabla \ref{tab:mapas_grandes_full} muestra los resultados de las pruebas con distintas combinaciones de $w$ y heurísticas.

\subsection{Pregunta D}
En mapas grandes y con alta densidad de obstáculos, AW-A* ofrece beneficios claros sobre A* y Greedy Best-First. 
Desafortunadamente nuestra implementacion de Greedy Best-First no encontro resultados, A* tarda demasiado en encontrar soluciones, AW-A* balancea eficiencia y calidad. 

\subsection{Pregunta E}
Un sistema anytime es preferible en escenarios donde el tiempo de respuesta es crítico, pero se dispone de tiempo extra para mejorar la solución. 
Ejemplos concretos incluyen planificación de rutas para drones, vehículos autónomos en entornos urbanos, o NPCs en videojuegos.

\subsection{Pregunta F}
Entre las limitaciones observadas destacan: sensibilidad a los parámetros ($w_{start}, w_{step}$), 
tiempos elevados en mapas muy grandes, y variabilidad en la calidad de la primera solución en mapas muy abiertos.

\section{Conclusiones}
AW-A* demuestra ser una estrategia intermedia entre la rapidez de Greedy y la optimalidad de A*. 
Permite obtener soluciones iniciales rápidas y mejorarlas progresivamente, lo que lo convierte en un enfoque adecuado para sistemas en tiempo real. 
Un caso adicional donde podría aplicarse es en la \textbf{logística de última milla}, donde cada segundo disponible puede refinar rutas de entrega.

\end{document}
